% TOP_PROBLEM: {"id": 30006677, "title": "Extended Riemann hypothesis for Dedekind zeta functions", "category_id": 1, "category": {"id": 1, "name": "number_theory", "display_name": "Number Theory", "description": "Properties of integers, prime numbers, Diophantine equations.", "slug": "number-theory", "order_index": 1, "created_at": "2026-07-31T15:26:25.670Z"}, "status": "open"}
% FIRST_POSTED: 2026-09-27
% !TeX program = lualatex
% ATTEMPT_STATUS: unresolved
% Research continuation by GPT_6_Astra_Ultra, 27 September 2026.
% Catalogue statement and existing sources preserved verbatim.
\documentclass[11pt]{article}
\usepackage[margin=25mm]{geometry}
\usepackage{fontspec}
\setmainfont{Latin Modern Roman}
\usepackage{amsmath,amssymb,mathtools}
\usepackage{unicode-math}
\setmathfont{Latin Modern Math}
\usepackage{xurl,hyperref}
\hypersetup{colorlinks=true,urlcolor=blue}
\setcounter{secnumdepth}{0}
\setlength{\parindent}{0pt}
\setlength{\parskip}{5pt}
\setlength{\emergencystretch}{3em}
\begin{document}
\section*{\texorpdfstring{Extended Riemann hypothesis for Dedekind zeta functions}{Extended Riemann hypothesis for Dedekind zeta functions}}\label{30006677}
\subsection{Definitions and mathematical statement}
For a number field $K$, let $\mathcal O_K$ be its ring of integers and $N\mathfrak a=\#(\mathcal O_K/\mathfrak a)$ the norm of a nonzero integral ideal. Define
\[
 \zeta_K(s)=\sum_{0\ne\mathfrak a\subseteq\mathcal O_K}(N\mathfrak a)^{-s}
 \qquad(\Re s>1),
\]
and continue meromorphically. The assertion is $\zeta_K(\rho)=0$ and $0<\Re\rho<1\Rightarrow\Re\rho=1/2$, for every finite extension $K/{\mathbb{Q}}$. The rational field itself is included.
\par\smallskip\textit{Formulation and concept references:} \href{https://www.aimath.org/WWN/rh/articles/html/18a/}{[S1]}.

\subsection{Short English statement}
For every number field, must the zeros of its ideal-counting zeta function inside the critical strip lie halfway across it?
\subsection{Research attempt}

\paragraph{Scope and checked background.}
Fix a finite extension $K/\mathbb Q$, of degree $d$, with absolute
discriminant $\Delta_K$. All constants in the argument may depend on this
field. The target subsequently quantifies over every such $K$; checking
quadratic or abelian fields alone is insufficient. Sutherland's notes
[E1], Remark 19.13, give the continuation, simple pole at $1$, and
functional equation of the completed Dedekind zeta function. Theorem 19.15
gives the factorization into primitive Dirichlet $L$-functions for abelian
extensions. These standard facts will be used with their stated scope.
The claims retained in catalogue sources [S3] and [S4] are not assumptions
of this attempt and have not been independently certified here.

The chosen route studies a once-integrated prime-ideal count. This smoothing
has a precise Mellin transform, so it permits a zero-exclusion argument
without an unjustified interchange over zeros. Greni\'e--Molteni [E2]
prove explicit estimates for this same type of smoothed count \emph{under}
GRH. Their work supplies established context for the scale being targeted,
not an unconditional bound that can be substituted into the missing step.
The work below derives the sufficient implication directly and examines
whether the elementary positivity of ideal coefficients can supply it.

\paragraph{Prime-power coefficients, including ramification.}
To avoid confusing ideals and integers, define a coefficient on positive
integers by
\[
 b_K(n)=\sum_{\substack{\mathfrak p,\,m\ge1\\(N\mathfrak p)^m=n}}
                 \log N\mathfrak p,
 \qquad \psi_K(x)=\sum_{n\le x}b_K(n).
\]
Unique factorization of ideals gives
\[
 -\frac{\zeta_K'}{\zeta_K}(s)=\sum_{n\ge1}b_K(n)n^{-s}
                   \qquad(\Re s>1).
\]
If $p\mathcal O_K=\prod_{j=1}^g\mathfrak p_j^{e_j}$ and
$N\mathfrak p_j=p^{f_j}$, then
\begin{equation}\label{ra:10:local}
 b_K(p^k)=\log p\sum_{j:\ f_j\mid k}f_j,
 \qquad 0\le b_K(p^k)\le d\log p.
\end{equation}
Here $\sum_j e_jf_j=d$ and $e_j\ge1$. Non-prime-powers have coefficient
zero. The exponents $e_j$ do not multiply the local logarithmic-derivative
coefficient; they enter only the degree constraint. Mistaking this point
would give incorrect contributions at ramified primes.

For completeness, absolute convergence on $\Re s>1$ follows from
\eqref{ra:10:local} and
$\sum_{n\ge2}(\log n)n^{-u}<\infty$ for $u>1$.
Thus differentiation of the Euler product and the integrations below
are legitimate in their initial domain. The positivity is genuine, but it
only compares these coefficients to $d$ times the ordinary von Mangoldt
coefficients. It does not yet identify the correct main term.

\paragraph{A smoothed criterion with the pole subtracted exactly.}
Define, for real $x\ge1$,
\[
 A_K(x)=\sum_{n\le x}b_K(n)(x-n)=\int_1^x\psi_K(t)\,dt,
 \qquad E_K(x)=A_K(x)-\frac{(x-1)^2}{2}.
\]
The subtraction is chosen for an exact identity. It differs from the usual
$x^2/2$ by only a linear polynomial, which is smaller than the contemplated
error. For $\Re s>1$, direct integration gives
\[
 \int_n^\infty(x-n)x^{-s-2}\,dx=\frac{n^{-s}}{s(s+1)},\qquad
 \int_1^\infty\frac{(x-1)^2}{2}x^{-s-2}\,dx
       =\frac1{(s-1)s(s+1)}.
\]
The first identity can be summed absolutely by the coefficient estimate.
Consequently
\begin{equation}\label{ra:10:mellin}
 -\frac{\zeta_K'}{\zeta_K}(s)
       =\frac1{s-1}+s(s+1)\int_1^\infty E_K(x)x^{-s-2}\,dx.
\end{equation}
The residue of $-\zeta_K'/\zeta_K$ at $1$ is one, regardless of the
field's degree or of the positive residue of $\zeta_K$ itself. This
explains why the main term is $x^2/2$, not $dx^2/2$.

\textbf{Lemma 10.1.} Let $0<\theta<1$. If
$E_K(x)=O_K(x^{1+\theta})$, then $\zeta_K$ has no zero in
$\Re s>\theta$.

\textit{Proof.}
The integral on the right of \eqref{ra:10:mellin} converges absolutely
and locally uniformly for $\Re s>\theta$. Its tails on a compact set
are bounded by a constant times
$\int_X^\infty x^{\theta-\Re s-1}\,dx$, tending uniformly to zero.
It therefore defines a holomorphic function there. The right side is
meromorphic with only the displayed pole at $1$, and equals the left side
initially for $\Re s>1$. The identity theorem extends the equality through
the half-plane. If $\rho\ne1$ were a zero of multiplicity $m\ge1$, then
writing $\zeta_K(s)=(s-\rho)^mh(s)$ with $h(\rho)\ne0$ would give
\[
 -\zeta_K'(s)/\zeta_K(s)=-\frac m{s-\rho}-h'(s)/h(s),
\]
a pole absent from the right side. This contradiction proves the lemma.
\hfill$\square$

It follows that the explicit target
\begin{equation}\label{ra:10:missing}
 \forall K\ \forall\varepsilon>0,\qquad
 E_K(x)=O_{K,\varepsilon}(x^{3/2+\varepsilon})
\end{equation}
would prove the catalogue assertion. Apply the lemma with
$\theta=1/2+\varepsilon<1$ and then let $\varepsilon$ range over positive
numbers. The completed functional equation reflects any remaining
left-hand strip zero to the right. Its gamma factors have no zeros or
poles in the open strip. This argument is independent of zero simplicity.
It also shows that a hypothetical zero of real part $\beta>1/2$ prevents
every bound $E_K=O(x^{1+\theta})$ with $\theta<\beta$; there is no need
to guess whether other oscillating zeros might cancel its contribution.

\paragraph{Attempt to obtain the estimate from local positivity.}
Equation \eqref{ra:10:local} gives the unconditional comparison
\[
 0\le A_K(x)\le dA_{\mathbb Q}(x).
\]
Even assuming an ideal estimate for the rational field, this comparison
allows an error of order $x^2$ after subtracting $x^2/2$ when $d>1$.
It therefore loses the degree-independent normalization that the pole
requires. The missing information is the distribution of splitting types
as $p$ varies, with sufficient cancellation at the square-root scale.
Replacing $b_K(p^k)$ everywhere by its largest permitted value discards
exactly that information.

The Gaussian field makes this failure explicit without any conjectural
factorization. A prime $p\equiv1\pmod4$ splits into two norm-$p$ primes;
a prime $p\equiv3\pmod4$ is inert with one norm-$p^2$ prime; $2$ ramifies
with one norm-$2$ prime. These assertions also follow by factoring
$X^2+1$ over $\mathbb F_p$ and checking $2=-i(1+i)^2$.
Let $\chi_4$ be the real character modulo four. Formula
\eqref{ra:10:local} becomes
\[
 b_{\mathbb Q(i)}(p^k)=(1+\chi_4(p^k))\log p.
\]
Thus
\[
 A_{\mathbb Q(i)}(x)=A_{\mathbb Q}(x)+B_4(x),\qquad
 B_4(x)=\sum_{p^k\le x}\chi_4(p^k)\log p\,(x-p^k).
\]
The desired estimate now needs control of a signed prime-power sum, not
just of the positive count. The parity of $k$ matters: at an inert prime,
odd powers contribute zero to $b_K$ and even powers contribute $2\log p$.
Counting only rational primes or only split primes would miss these terms.

This same identity proves the Euler-factor relation
$\zeta_{\mathbb Q(i)}=\zeta L(s,\chi_4)$ in the absolutely convergent
region, hence meromorphically everywhere. In the open strip both factors
are holomorphic, so their zero orders add. There can be no cancellation
of a zero by a pole there. In particular, proving RH for $\zeta$ alone
does not settle this first quadratic example: a zero of the beta function
$L(s,\chi_4)$ would still be a Dedekind zeta zero.

\paragraph{Abelian reduction and its exact stopping point.}
More generally, the abelian factorization in [E1] shows that Dirichlet GRH
implies the Dedekind assertion for every finite abelian extension of
$\mathbb Q$. Conversely, applying the Dedekind assertion to cyclotomic
fields covers all primitive Dirichlet characters, since each occurs as
a holomorphic factor of the appropriate cyclotomic zeta function in the
strip. These implications are about the abelian subfamily. A general
number field need not be contained in an abelian extension of $\mathbb Q$,
so this does not reduce the full catalogue target to Dirichlet characters.
No unproved holomorphy assertion for general Artin $L$-functions is being
inserted to fill that gap.

The smoothing itself also needs a stress test. For $\beta\in(1/2,1)$
and $\gamma\ne0$, the elementary model
$F(x)=\Re(x^{1+\beta+i\gamma})$ has Mellin transform
\[
 \int_1^\infty F(x)x^{-s-2}\,dx
   =\frac12\left(\frac1{s-\beta-i\gamma}
                        +\frac1{s-\beta+i\gamma}\right)
       \qquad(\Re s>\beta).
\]
Positive and negative oscillations do not remove either pole. The factor
$s(s+1)$ in \eqref{ra:10:mellin} is nonzero at both points. This toy
calculation is not an explicit formula for an actual field; it verifies
why integration improves convergence without erasing a possible off-line
obstruction. The meromorphic argument in Lemma 10.1 supplies the rigorous
statement for the actual zeta function.

\paragraph{Would additional smoothing close the gap?}
For any integer $m\ge1$, consider
$A_{K,m}(x)=\sum_{n\le x}b_K(n)(x-n)^m/m!$.
Successive integration, or the substitution $x=n/u$, gives
\[
 \int_n^\infty\frac{(x-n)^m}{m!}x^{-s-m-1}\,dx
        =\frac{n^{-s}}{s(s+1)\cdots(s+m)}\qquad(\Re s>0).
\]
In the initial half-plane $\Re s>1$ this can again be summed absolutely.
Subtracting $(x-1)^{m+1}/(m+1)!$ removes exactly the pole term
$1/(s-1)$ after multiplying the transform by
$s(s+1)\cdots(s+m)$. Therefore an error
$O(x^{m+1/2+\varepsilon})$ would still exclude the same off-line zeros.
The proof is the identical holomorphic-continuation argument, with a
different polynomial denominator. None of its factors vanishes in the
open critical strip.

This calculation locates the limit of the smoothing strategy. Integrating
an available bound $E_K(x)=O(x^{1+\theta})$ raises its exponent by one;
it does not decrease $\theta$. Positivity permits Fubini in a larger
collection of finite calculations, but supplies no reduction of the
relative error exponent. In particular, many integrations of a coarse
degree-based majorant cannot turn it into the desired square-root
discrepancy. A genuine new estimate must enter before the Mellin
continuation step. This also explains why the conditional estimates in
[E2] cannot be reversed into an unconditional proof merely by applying
an inverse transform: their zero-location hypothesis already supplies
the cancellation needed to control those transforms.

\paragraph{Checks and what they verify.}
The local Python script checks the Gaussian coefficient identity exactly
at every one of the $1280$ prime powers at most $10^4$, including all powers
of $2$. It evaluates the finite smoothed sums using ordinary floating-point
logarithms, obtaining
\[
\begin{array}{r|rr}
x&E_{\mathbb Q}(x)/x^{3/2}&E_{\mathbb Q(i)}(x)/x^{3/2}\\\hline
10&-0.213006&-0.351583\\
100&-0.095511&-0.244875\\
1000&-0.036110&-0.092366\\
10000&-0.007322&0.059332
\end{array}
\]
The change of sign in the last column is consistent with a signed error;
no sign assertion about $E_K$ was used. These decimal values are diagnostics,
not certified inequalities, and the four points cannot establish a
big-$O$ estimate. Exact rational checks at three real $s>1$ separately
verify the subtraction identity preceding \eqref{ra:10:mellin}.
The source and output are stored in
\texttt{\_Local/top\_problems/continuation\_2026\_09\_27/analytic/}
as \texttt{checks.py} and \texttt{results.json}.

\paragraph{First gap and outcome.}
All implications from \eqref{ra:10:missing} to the desired zero location
are justified, including multiplicities, the pole at one, gamma factors,
ramification, and field dependence of constants. The first missing step
is the estimate \eqref{ra:10:missing} itself for each fixed field.
Local positivity proves only a degree-dependent majorant and gives no
square-root cancellation between splitting types. Even a successful
argument for all abelian fields would leave the nonabelian scope untouched.
An effective next lemma would have to bound the signed discrepancy of
the local coefficients from their pole-normalized mean, uniformly in $x$
for the chosen field; the Gaussian formula gives the first concrete test.

\textbf{Outcome: UNRESOLVED.} The checked output is a direct smoothed
Mellin reduction and a local-to-global obstruction analysis, with finite
arithmetic diagnostics. These are not claimed as new theorems establishing
ERH. No unconditional square-root error estimate and no off-line zero have
been obtained. Primary-source inspection and explicit calculations were
used; independent expert verification was not performed.

\subsection{Sources}
\begin{itemize}
\item[S1]American Institute of Mathematics RH project, The Extended Riemann Hypothesis; undated article18a.
\url{https://www.aimath.org/WWN/rh/articles/html/18a/}.
\textit{Formulation source.}
\item[S2]A random polynomial with multiplicative coefficients is almost surely irreducible — Varjú and Xu.
\url{https://academic.oup.com/imrn/article/2026/16/rnag178/8761102}.
\textit{Further reference; not a proof endorsement.}
\item[S3]On the Extended, Generalized, and Grand Riemann Hypotheses: A Unified Approach Based on the General Properties of L-Functions — Weicun Zhang, v19.
\url{https://www.preprints.org/manuscript/202506.0481}.
\textit{Further reference; not a proof endorsement.}
\item[S4]Proof of generalized Riemann hypothesis for Dedekind zetas and Dirichlet L-functions — Andrzej Mądrecki.
\url{https://arxiv.org/abs/0706.0256}.
\textit{Further reference; not a proof endorsement.}

\item[E1]Andrew V. Sutherland, \emph{18.785 Number Theory I, Lecture 19:
The analytic class number formula}, MIT, fall 2021,
Remark 19.13 and Theorem 19.15.
\url{https://math.mit.edu/classes/18.785/2021fa/LectureNotes19.pdf}.
\textit{Author lecture notes: completed functional equation and factorization
for abelian extensions. Consulted 27 September 2026.}
\item[E2]Lo\"ic Greni\'e and Giuseppe Molteni,
\emph{Explicit smoothed prime ideals theorems under GRH},
Mathematics of Computation 85 (2016), no. 300, 1875--1899,
DOI: 10.1090/mcom3039; arXiv:1312.4465.
\url{https://arxiv.org/abs/1312.4465}.
Author manuscript:
\url{https://sites.unimi.it/molteni/research/papers-pdf/34-molteni-Explicit_smoothed_prime_ideals_theorems_under_GRH.pdf}.
\textit{Primary research context: estimates explicitly assume GRH;
no conditional estimate is used unconditionally here.
Consulted 27 September 2026.}
\end{itemize}
\end{document}
